\chapter{Inleiding}

In dit verslag wordt het ontwerp, de productie, de berekening en de beproeving
van een glasvezelversterkte epoxy trektestcoupon uitgewerkt. Het doel is om de
verwachte sterkte en stijfheid van de coupon te bepalen, deze waarden te
vergelijken met een ELaminate-simulatie en de resultaten daarna naast de
werkelijke trekproef te leggen.

Een vezelversterkte kunststof bestaat uit twee hoofdbestanddelen: vezels en een
matrix. De glasvezels nemen vooral de trekbelasting op en zorgen voor de hoge
stijfheid en sterkte in vezelrichting. De epoxy matrix houdt de vezels op hun
plaats, verdeelt de belasting tussen de vezels en beschermt het laminaat tegen
schade van buitenaf. De eigenschappen van het composiet worden daarom bepaald
door zowel de materiaalkeuze als de vezelrichting, het vezelvolumepercentage,
de hechting tussen vezel en matrix en de kwaliteit van het productieproces.

Een afzonderlijke laag in het composiet wordt een lamel genoemd. Meerdere
lamellen samen vormen een laminaat. In dit verslag bestaat het laminaat uit vier
lamellen met de opbouw:

\[
[-30^\circ / 45^\circ / -45^\circ / 30^\circ]
\]

Daarbij is de globale \(x\)-richting gelijk aan de trekrichting van de coupon.
Elke lamel heeft ook een lokaal assenstelsel: richting 1 ligt in de
vezelrichting en richting 2 ligt dwars op de vezelrichting. De component 12
hoort bij de schuifvervorming in het vlak van de lamel. Omdat de vezels niet in de
trekrichting liggen, moeten de lokale lamel-eigenschappen worden omgerekend
naar het globale assenstelsel van de coupon.

Voor deze omrekening wordt de klassieke laminaattheorie gebruikt. Eerst worden
de equivalente lamel-eigenschappen bepaald met mengregels. Daarna wordt de
lokale stijfheidsmatrix \(Q\) opgesteld en per vezelhoek getransformeerd naar
\(\bar{Q}\). Met deze getransformeerde matrices worden de \(A\)-, \(B\)- en
\(D\)-matrix bepaald. Daarmee kunnen rekken, spanningen, torsievervorming,
stijfheid en een geschatte faalkracht worden berekend.

De productie van de coupon is belangrijk, omdat composietmateriaal tijdens het
maken pas zijn uiteindelijke eigenschappen krijgt. Luchtinsluitingen, een
ongelijke hoeveelheid hars, verschoven vezelrichtingen, onvoldoende vacuüm of
inklemschade tijdens de trekproef kunnen ervoor zorgen dat de gemeten waarden
afwijken van de berekende waarden. Daarom wordt in het verslag niet alleen de
berekening uitgewerkt, maar ook het productieproces, de afwijkingen tijdens de
productie, de trekproef en het waargenomen faalmechanisme.

De opbouw van het verslag volgt deze lijn. Eerst worden de coupon, norm,
materiaaleigenschappen en het productievoorschrift uitgewerkt. Daarna wordt het
werkelijke productieproces beschreven. Vervolgens wordt de coupon berekend met
de klassieke laminaattheorie en gesimuleerd met ELaminate. Tot slot worden de
trekproefresultaten terugvertaald naar lokale lamelspanningen en vergeleken met
de berekende sterkte en stijfheid.
